\subsection{Relation to previous work}
\label{subsec:related-work}

The tournament case of Seymour's second neighborhood conjecture was proved
by Fisher~\cite{Fisher1996} using a probability distribution obtained from
a skew-symmetric game. Havet and Thomass\'e~\cite{HavetThomasse2000}
subsequently gave a proof using median orders, including a stronger bound
involving the good vertices of an order. Both methods enter the present
work. The incoming-core argument uses the existence of a losing density
on a tournament; the completion arguments use the strong feed bound and
the equality case of sedimentation. These are established inputs. Our
arguments specify the additional conditions needed to transfer information
from a tournament completion to the exact second neighborhoods of the
original graph. In particular, the existence of a Seymour vertex in a
completion alone does not provide such a transfer.

Fidler and Yuster~\cite{FidlerYuster2007} proved the vertex-weighted
tournament statement and established further cases obtained by deleting
prescribed edge sets from tournaments. Ghazal's work on generalized stars
and other missing graphs~\cite{Ghazal2012,Ghazal2013} develops the convenient
orientation and good-missing-edge approach underlying several of our
completion definitions. Dara, Francis, Jacob and
Narayanan~\cite{DaraEtAl2022} extended the matching case to a missing graph
whose edges can be partitioned into a matching and a star, and also
considered partitions into an independent set and a $2$-degenerate graph.
Our protection-preserving completions require additional compatibility
conditions beyond convenience of individual added arcs. The proof of each
transfer inequality therefore accounts for paths using added arcs and for
original second neighbors that become first neighbors. None of these
inequalities is inferred simply by applying a known tournament theorem to
an arbitrary completion.

Minimal-counterexample methods provide another starting point.
Brantner, Brockman, Kay and Snively~\cite{BrantnerEtAl2009} established
strong connectivity, gaps in $\{1,2\}$, and restrictions involving
gap-one predecessors and directed cycles for their minimal counterexamples.
Their extremal selection first minimizes the number of arcs over all
counterexamples and then minimizes the order. The selections used here
must be distinguished from that convention: fixed-order minimality under
deletion of a nonempty arc set, and minimum order followed by arc count
and further finite objectives, serve different arguments. We give the
required deletion and copying arguments for the stated selection in each
case. In particular, the predecessor-containment row operation and the
resulting normal form are not treated as consequences of the earlier
extremal definition.

Seacrest's subset formulation~\cite{SeacrestSubsets2019} supplies a closely
related boundary perspective. Its edge-minimal subset inequality motivates
the boundary budgets used in the residual analysis. Because the sets and
minimality assumptions have to agree exactly, the external-boundary form
needed here is proved explicitly. This paper is distinct from Seacrest's
arc-weighted extension~\cite{Seacrest2015}; an arc-weighted result is not
substituted for the unweighted feed or boundary statements. The residual
bound $R\geq 7$ in this manuscript concerns the specified minimal
all-deficient graphs. It must not be read as an assertion about the
missing-edge count of every counterexample before a minimizing deletion.

There are further reasons to keep extremal assumptions visible.
Espuny D\'iaz, Gir\~ao, Granet and
Kronenberg~\cite{EspunyDiazEtAl2025} prove that a vertex-minimal
counterexample has minimum outdegree greater than the square root of its
order, while also constructing arbitrarily large strongly connected
counterexamples with bounded minimum outdegree, conditional on the
existence of any counterexample. Zelenskiy, Darmosiuk and
Nalivayko~\cite{ZelenskiyEtAl2021} obtain conditional counterexample
constructions with widely differing densities and diameters. Thus
counterexample amplification, and cyclic substitution in particular,
already have a literature. The calculations here involving such
constructions are used to delimit particular completion or weighted
transfer claims, rather than to assert a new general amplification
principle.

Theorem~\ref{thm:core-five} has a different parameter from these density
and degree results: the number of distinct inclusion-maximal incoming
neighborhoods. There is no restriction on the order of the original graph
or on the size of an incoming-neighborhood class. The certificate on a
chosen core is checked against every permitted subset of an incoming set,
so that it controls all original target patterns simultaneously. An
ordinary verification of Seymour's conjecture for small graphs does not
automatically supply these stronger certificates. The finite tables are
therefore accompanied by explicit integer checks and a proof of their
transfer to arbitrary-order graphs. The labelled construction from an
infeasible pattern system establishes a converse at the universal level.
Finite linear-programming duality is standard; the graph realization and
its exact-distance calculation are the substantive steps that must be
verified for this particular reformulation. Universal feasibility remains
unproved.

Recent special cases provide useful comparisons. Wang and
Lu~\cite{WangLu2026} prove the conjecture for
$\mathcal D_{0,3}\cup\mathcal D_{1,1}$, where $\mathcal D_{s,t}$ consists
of oriented graphs admitting a vertex partition whose induced underlying
graphs are respectively $s$- and $t$-degenerate. Bai, Li and
Park~\cite{BaiLiPark2026} study the stronger requirement of a complete
directed matching from the first to the exact second neighborhood, proving
it when the minimum outdegree is at most five or the graph is
$5$-anti-transitive. Kaneko and Locke~\cite{KanekoLocke2001} established
ordinary Seymour vertices when the minimum outdegree is at most six.
The preprint of Sadhukhan, Sandeep and Sen~\cite{SadhukhanEtAl2026}
treats minimum outdegree seven using computer-assisted local reductions.
Brukhman's preprint~\cite{Brukhman2026} treats the dense range
$n\leq 2\delta^++2$. The incoming-core theorem uses none of these
recent preprints as a premise.

The five-class hypothesis does not imply the hypotheses of those
particular recent results. For example, replace each vertex of the regular
five-vertex tournament, with arcs $i\to i+1,i+2$ modulo five, by an
independent set of size $m\geq4$. The resulting graph has five maximal
incoming classes, order $5m$, and minimum outdegree $2m$. Its underlying
complete five-partite graph belongs to neither
$\mathcal D_{0,3}$ nor $\mathcal D_{1,1}$, and a path through parts
$0,1,2,3,4,1$ excludes $5$-anti-transitivity. It also lies outside the
stated dense range. This example itself is already covered by the weighted
tournament theorem. It demonstrates noncontainment in specific hypotheses,
not novelty relative to every known weighted or substitution argument.

Halkiewicz's split-twin operation~\cite{Halkiewicz2026} duplicates incoming
neighbors and partitions outgoing neighbors while increasing the order;
its stated preservation theorem has a condition involving an existing
Seymour vertex. This differs from the fixed-order row copying used here.
Nevertheless, this distinction and the preceding comparisons do not settle
priority for equivalent formulations under other terminology. The present
paper makes no claim to an exhaustive classification of earlier reductions.
The literature comparison is current to 7 October 2026. Public full-proof
claims, including~\cite{GloverClaim2026}, are not assumed correct or used
as premises. The conclusions proved here remain partial results and
reformulations of the conjecture.
